% Generated draft: requires mathematical and source review.
\documentclass[11pt]{article}
\usepackage[margin=25mm]{geometry}
\usepackage{fontspec}
\setmainfont{Latin Modern Roman}
\usepackage{amsmath,amssymb,mathtools,unicode-math}
\setmathfont{Latin Modern Math}
\usepackage{enumitem,needspace,xurl,hyperref}
\setlength{\parindent}{0pt}
\setlength{\parskip}{5pt}
\setlength{\emergencystretch}{3em}
\newcommand{\N}{\mathbb{N}}
\newcommand{\Z}{\mathbb{Z}}
\newcommand{\Q}{\mathbb{Q}}
\newcommand{\R}{\mathbb{R}}
\newcommand{\C}{\mathbb{C}}
\newcommand{\F}{\mathbb{F}}
\newcommand{\A}{\mathbb{A}}
\newcommand{\PP}{\mathbb{P}}
\newcommand{\E}{\mathbb{E}}
\newcommand{\eps}{\varepsilon}
\newcommand{\dd}{\,\mathrm{d}}
\begin{document}
\section*{UnsolvedMath ID: 30006214}
\textbf{Source problem code:} catalogue entry
\subsection{Definitions and mathematical statement}
\textbf{Source question.} Find an intrinsic incidence-variety description of the Euler discriminant, analogous to the geometric description intended by the source for classical discriminants.

\textbf{Meaning and unresolved scope.} An incidence-variety description would specify an incidence variety \(I\), a parameter space \(S\), and a projection \(\pi\colon I\to S\) from which the Euler discriminant can be recovered. At the level of underlying closed sets, a possible schematic form of such a description is
\[
\lvert\mathcal D_{\mathrm{Euler}}\rvert=\overline{\pi(I)}.
\]
This display is only a template explaining the requested type of description; it is not a definition or an asserted identity supplied by the record. The record does not specify the input geometric family, the parameter space, the ground field, the ambient space for the incidence variety, or the definition of the Euler discriminant. It also does not specify whether recovery means equality of sets, schemes, divisors, or cycles, or whether taking a closure is appropriate. Consequently, no definite quantification over geometric families or precise incidence conditions can be extracted without additional source context.

The original question asks whether one can find the same picture for the Euler discriminant, but the antecedent classical-discriminant picture is not included in the supplied record. Thus the required analogy and the meaning of intrinsic remain unresolved.

\textbf{Supplied status qualification.} The record labels the problem partially solved. Its supplied literature assessment states that Euler-discriminant elimination is proved for important families, including hyperplane complements, while a general intrinsic incidence-variety description remains unresolved. These are supplied dated curation claims, not independently verified mathematical or literature conclusions.
\subsection{Short English statement}
The problem asks for a geometric incidence construction that intrinsically recovers the Euler discriminant, analogous to a classical-discriminant construction. The supplied record omits the defining family, discriminant convention, and comparison construction, so the intended general statement cannot yet be made fully precise. The supplied partially solved status requires independent review.
\subsection{Sources}
\begin{itemize}
\item Local UnsolvedMath registry, supplied record with permanent ID \(30006214\), source problem code \texttt{catalogue entry}, title ``Incident-Variety Description of Euler Discriminants.'' The project identifies \texttt{data/unsolvedmath/problems.json} and the supplementary registry \texttt{catalog/problems.json}; the supplied record's location within these files was not independently checked.
\item Supplied report citation: \emph{MFO-RIMS Tandem Workshop: Optimization, Theoretical Computer Science and Algebraic Geometry: Convexity and Beyond} (2025), Oberwolfach Reports, pp.\ 405--450, DOI 10.4171/owr/2025/9. \url{https://doi.org/10.4171/owr/2025/9}. The record attributes the question to the workshop report and gives proposed year 2025, but names no individual proposer and supplies no precise passage locator.
\item Supplied literature source: \emph{Euler discriminants of hyperplane complements}. \url{https://arxiv.org/abs/2411.19696}.
\item Supplied literature source: \emph{Hypergeometric discriminants}. \url{https://arxiv.org/abs/2505.13163}.
\item Literature assessment date: \texttt{2026-08-22}. The partially solved label, the claim of elimination results for important families including hyperplane complements, and the claim that the intrinsic general description remains unresolved are supplied dated curation derived from the OWR material, not independently verified. The supplied background explicitly requests independent verification of time-sensitive status and cited claims. None of the linked sources was visited for this draft.
\end{itemize}
\end{document}
